% Options for packages loaded elsewhere
% Options for packages loaded elsewhere
\PassOptionsToPackage{unicode}{hyperref}
\PassOptionsToPackage{hyphens}{url}
\PassOptionsToPackage{dvipsnames,svgnames,x11names}{xcolor}
%
\documentclass[
  10pt,
  a4paper,10pt]{extarticle}
\usepackage{xcolor}
\usepackage[lmargin=1in,rmargin=1in,tmargin=0.8in,bmargin=0.8in]{geometry}
\usepackage{amsmath,amssymb}
\setcounter{secnumdepth}{5}
\usepackage{iftex}
\ifPDFTeX
  \usepackage[T1]{fontenc}
  \usepackage[utf8]{inputenc}
  \usepackage{textcomp} % provide euro and other symbols
\else % if luatex or xetex
  \usepackage{unicode-math} % this also loads fontspec
  \defaultfontfeatures{Scale=MatchLowercase}
  \defaultfontfeatures[\rmfamily]{Ligatures=TeX,Scale=1}
\fi
\usepackage{lmodern}
\ifPDFTeX\else
  % xetex/luatex font selection
\fi
% Use upquote if available, for straight quotes in verbatim environments
\IfFileExists{upquote.sty}{\usepackage{upquote}}{}
\IfFileExists{microtype.sty}{% use microtype if available
  \usepackage[]{microtype}
  \UseMicrotypeSet[protrusion]{basicmath} % disable protrusion for tt fonts
}{}
\makeatletter
\@ifundefined{KOMAClassName}{% if non-KOMA class
  \IfFileExists{parskip.sty}{%
    \usepackage{parskip}
  }{% else
    \setlength{\parindent}{0pt}
    \setlength{\parskip}{6pt plus 2pt minus 1pt}}
}{% if KOMA class
  \KOMAoptions{parskip=half}}
\makeatother
% Make \paragraph and \subparagraph free-standing
\makeatletter
\ifx\paragraph\undefined\else
  \let\oldparagraph\paragraph
  \renewcommand{\paragraph}{
    \@ifstar
      \xxxParagraphStar
      \xxxParagraphNoStar
  }
  \newcommand{\xxxParagraphStar}[1]{\oldparagraph*{#1}\mbox{}}
  \newcommand{\xxxParagraphNoStar}[1]{\oldparagraph{#1}\mbox{}}
\fi
\ifx\subparagraph\undefined\else
  \let\oldsubparagraph\subparagraph
  \renewcommand{\subparagraph}{
    \@ifstar
      \xxxSubParagraphStar
      \xxxSubParagraphNoStar
  }
  \newcommand{\xxxSubParagraphStar}[1]{\oldsubparagraph*{#1}\mbox{}}
  \newcommand{\xxxSubParagraphNoStar}[1]{\oldsubparagraph{#1}\mbox{}}
\fi
\makeatother

\usepackage{color}
\usepackage{fancyvrb}
\newcommand{\VerbBar}{|}
\newcommand{\VERB}{\Verb[commandchars=\\\{\}]}
\DefineVerbatimEnvironment{Highlighting}{Verbatim}{commandchars=\\\{\}}
% Add ',fontsize=\small' for more characters per line
\usepackage{framed}
\definecolor{shadecolor}{RGB}{241,243,245}
\newenvironment{Shaded}{\begin{snugshade}}{\end{snugshade}}
\newcommand{\AlertTok}[1]{\textcolor[rgb]{0.68,0.00,0.00}{#1}}
\newcommand{\AnnotationTok}[1]{\textcolor[rgb]{0.37,0.37,0.37}{#1}}
\newcommand{\AttributeTok}[1]{\textcolor[rgb]{0.40,0.45,0.13}{#1}}
\newcommand{\BaseNTok}[1]{\textcolor[rgb]{0.68,0.00,0.00}{#1}}
\newcommand{\BuiltInTok}[1]{\textcolor[rgb]{0.00,0.23,0.31}{#1}}
\newcommand{\CharTok}[1]{\textcolor[rgb]{0.13,0.47,0.30}{#1}}
\newcommand{\CommentTok}[1]{\textcolor[rgb]{0.37,0.37,0.37}{#1}}
\newcommand{\CommentVarTok}[1]{\textcolor[rgb]{0.37,0.37,0.37}{\textit{#1}}}
\newcommand{\ConstantTok}[1]{\textcolor[rgb]{0.56,0.35,0.01}{#1}}
\newcommand{\ControlFlowTok}[1]{\textcolor[rgb]{0.00,0.23,0.31}{\textbf{#1}}}
\newcommand{\DataTypeTok}[1]{\textcolor[rgb]{0.68,0.00,0.00}{#1}}
\newcommand{\DecValTok}[1]{\textcolor[rgb]{0.68,0.00,0.00}{#1}}
\newcommand{\DocumentationTok}[1]{\textcolor[rgb]{0.37,0.37,0.37}{\textit{#1}}}
\newcommand{\ErrorTok}[1]{\textcolor[rgb]{0.68,0.00,0.00}{#1}}
\newcommand{\ExtensionTok}[1]{\textcolor[rgb]{0.00,0.23,0.31}{#1}}
\newcommand{\FloatTok}[1]{\textcolor[rgb]{0.68,0.00,0.00}{#1}}
\newcommand{\FunctionTok}[1]{\textcolor[rgb]{0.28,0.35,0.67}{#1}}
\newcommand{\ImportTok}[1]{\textcolor[rgb]{0.00,0.46,0.62}{#1}}
\newcommand{\InformationTok}[1]{\textcolor[rgb]{0.37,0.37,0.37}{#1}}
\newcommand{\KeywordTok}[1]{\textcolor[rgb]{0.00,0.23,0.31}{\textbf{#1}}}
\newcommand{\NormalTok}[1]{\textcolor[rgb]{0.00,0.23,0.31}{#1}}
\newcommand{\OperatorTok}[1]{\textcolor[rgb]{0.37,0.37,0.37}{#1}}
\newcommand{\OtherTok}[1]{\textcolor[rgb]{0.00,0.23,0.31}{#1}}
\newcommand{\PreprocessorTok}[1]{\textcolor[rgb]{0.68,0.00,0.00}{#1}}
\newcommand{\RegionMarkerTok}[1]{\textcolor[rgb]{0.00,0.23,0.31}{#1}}
\newcommand{\SpecialCharTok}[1]{\textcolor[rgb]{0.37,0.37,0.37}{#1}}
\newcommand{\SpecialStringTok}[1]{\textcolor[rgb]{0.13,0.47,0.30}{#1}}
\newcommand{\StringTok}[1]{\textcolor[rgb]{0.13,0.47,0.30}{#1}}
\newcommand{\VariableTok}[1]{\textcolor[rgb]{0.07,0.07,0.07}{#1}}
\newcommand{\VerbatimStringTok}[1]{\textcolor[rgb]{0.13,0.47,0.30}{#1}}
\newcommand{\WarningTok}[1]{\textcolor[rgb]{0.37,0.37,0.37}{\textit{#1}}}

\usepackage{longtable,booktabs,array}
\usepackage{calc} % for calculating minipage widths
% Correct order of tables after \paragraph or \subparagraph
\usepackage{etoolbox}
\makeatletter
\patchcmd\longtable{\par}{\if@noskipsec\mbox{}\fi\par}{}{}
\makeatother
% Allow footnotes in longtable head/foot
\IfFileExists{footnotehyper.sty}{\usepackage{footnotehyper}}{\usepackage{footnote}}
\makesavenoteenv{longtable}
\usepackage{graphicx}
\makeatletter
\newsavebox\pandoc@box
\newcommand*\pandocbounded[1]{% scales image to fit in text height/width
  \sbox\pandoc@box{#1}%
  \Gscale@div\@tempa{\textheight}{\dimexpr\ht\pandoc@box+\dp\pandoc@box\relax}%
  \Gscale@div\@tempb{\linewidth}{\wd\pandoc@box}%
  \ifdim\@tempb\p@<\@tempa\p@\let\@tempa\@tempb\fi% select the smaller of both
  \ifdim\@tempa\p@<\p@\scalebox{\@tempa}{\usebox\pandoc@box}%
  \else\usebox{\pandoc@box}%
  \fi%
}
% Set default figure placement to htbp
\def\fps@figure{htbp}
\makeatother


% definitions for citeproc citations
\NewDocumentCommand\citeproctext{}{}
\NewDocumentCommand\citeproc{mm}{%
  \begingroup\def\citeproctext{#2}\cite{#1}\endgroup}
\makeatletter
 % allow citations to break across lines
 \let\@cite@ofmt\@firstofone
 % avoid brackets around text for \cite:
 \def\@biblabel#1{}
 \def\@cite#1#2{{#1\if@tempswa , #2\fi}}
\makeatother
\newlength{\cslhangindent}
\setlength{\cslhangindent}{1.5em}
\newlength{\csllabelwidth}
\setlength{\csllabelwidth}{3em}
\newenvironment{CSLReferences}[2] % #1 hanging-indent, #2 entry-spacing
 {\begin{list}{}{%
  \setlength{\itemindent}{0pt}
  \setlength{\leftmargin}{0pt}
  \setlength{\parsep}{0pt}
  % turn on hanging indent if param 1 is 1
  \ifodd #1
   \setlength{\leftmargin}{\cslhangindent}
   \setlength{\itemindent}{-1\cslhangindent}
  \fi
  % set entry spacing
  \setlength{\itemsep}{#2\baselineskip}}}
 {\end{list}}
\usepackage{calc}
\newcommand{\CSLBlock}[1]{\hfill\break\parbox[t]{\linewidth}{\strut\ignorespaces#1\strut}}
\newcommand{\CSLLeftMargin}[1]{\parbox[t]{\csllabelwidth}{\strut#1\strut}}
\newcommand{\CSLRightInline}[1]{\parbox[t]{\linewidth - \csllabelwidth}{\strut#1\strut}}
\newcommand{\CSLIndent}[1]{\hspace{\cslhangindent}#1}



\setlength{\emergencystretch}{3em} % prevent overfull lines

\providecommand{\tightlist}{%
  \setlength{\itemsep}{0pt}\setlength{\parskip}{0pt}}



 


\makeatletter
\@ifpackageloaded{caption}{}{\usepackage{caption}}
\AtBeginDocument{%
\ifdefined\contentsname
  \renewcommand*\contentsname{Table of contents}
\else
  \newcommand\contentsname{Table of contents}
\fi
\ifdefined\listfigurename
  \renewcommand*\listfigurename{List of Figures}
\else
  \newcommand\listfigurename{List of Figures}
\fi
\ifdefined\listtablename
  \renewcommand*\listtablename{List of Tables}
\else
  \newcommand\listtablename{List of Tables}
\fi
\ifdefined\figurename
  \renewcommand*\figurename{Figure}
\else
  \newcommand\figurename{Figure}
\fi
\ifdefined\tablename
  \renewcommand*\tablename{Table}
\else
  \newcommand\tablename{Table}
\fi
}
\@ifpackageloaded{float}{}{\usepackage{float}}
\floatstyle{ruled}
\@ifundefined{c@chapter}{\newfloat{codelisting}{h}{lop}}{\newfloat{codelisting}{h}{lop}[chapter]}
\floatname{codelisting}{Listing}
\newcommand*\listoflistings{\listof{codelisting}{List of Listings}}
\makeatother
\makeatletter
\makeatother
\makeatletter
\@ifpackageloaded{caption}{}{\usepackage{caption}}
\@ifpackageloaded{subcaption}{}{\usepackage{subcaption}}
\makeatother
\usepackage{bookmark}
\IfFileExists{xurl.sty}{\usepackage{xurl}}{} % add URL line breaks if available
\urlstyle{same}
\hypersetup{
  pdftitle={Design and analysis of experiments},
  pdfauthor={Macartan Humphreys},
  colorlinks=true,
  linkcolor={blue},
  filecolor={Maroon},
  citecolor={Blue},
  urlcolor={orange},
  pdfcreator={LaTeX via pandoc}}


\title{Design and analysis of experiments}
\usepackage{etoolbox}
\makeatletter
\providecommand{\subtitle}[1]{% add subtitle to \maketitle
  \apptocmd{\@title}{\par {\large #1 \par}}{}{}
}
\makeatother
\subtitle{HU \textbar{} Spring 2026 (Winter term 25-26)}
\author{Macartan Humphreys}
\date{}
\begin{document}
\maketitle
\begin{abstract}
The course addresses the design and analysis of experiments. It is a
hands-on course in which theoretical results are introduced through
lecture, demonstrated in practice, and worked on in groups. Topics
include randomization schemes, analysis strategies, design evaluation,
and workflows.
\end{abstract}


\section{Class resources}\label{class-resources}

\begin{itemize}
\item
  All resources: \url{https://macartan.github.io/ci}
\item
  This syllabus:
  \url{https://macartan.github.io/ci/experiments_syllabus_2026.pdf}
\item
  Slides: \url{https://macartan.github.io/ci/experiments_2026.html}
\item
  \href{https://cloud.wzb.eu/apps/forms/s/8QmokT5GeQfkmkrBHDzrzN4j}{Student
  survey}
\item
  Git repo \url{https://github.com/macartan/ci}: This repo is for both
  my causal inference class and this experiments class.
\end{itemize}

\section{Times and locations}\label{times-and-locations}

Meetings will be held at the WZB

\begin{longtable}[]{@{}
  >{\raggedright\arraybackslash}p{(\linewidth - 6\tabcolsep) * \real{0.1538}}
  >{\raggedright\arraybackslash}p{(\linewidth - 6\tabcolsep) * \real{0.1385}}
  >{\raggedright\arraybackslash}p{(\linewidth - 6\tabcolsep) * \real{0.2154}}
  >{\raggedright\arraybackslash}p{(\linewidth - 6\tabcolsep) * \real{0.4923}}@{}}
\toprule\noalign{}
\begin{minipage}[b]{\linewidth}\raggedright
Date
\end{minipage} & \begin{minipage}[b]{\linewidth}\raggedright
Time
\end{minipage} & \begin{minipage}[b]{\linewidth}\raggedright
Location
\end{minipage} & \begin{minipage}[b]{\linewidth}\raggedright
Topic
\end{minipage} \\
\midrule\noalign{}
\endhead
\bottomrule\noalign{}
\endlastfoot
9 Jan 2026 & 9 - 13 & WZB A 305 & Getting started \\
16 Jan 2026 & 9 - 13 & WZB A 305 & Doing experiments \\
30 Jan 2026 & 9 - 13 & WZB A 305 & Causality and causal inquiries \\
6 Feb 2026 & 9 - 13 & WZB A 305 & Answer strategies \\
20 Feb 2026 & 9 - 13 & WZB A 305 & Data strategies and design
evaluation \\
27 Feb 2026 & 9 - 13 & WZB B 001 & Topics: downstream experimentation,
spillover experiments, patient preference trials, survey experiments \\
\end{longtable}

The sessions will be structured roughly as follows.

\begin{itemize}
\tightlist
\item
  9:00 - 11:00: lecture or discussions of readings
\item
  11:00 - 11:15: break
\item
  11:15 - 13:00: discussions, team presentations and in-class exercises
\end{itemize}

\section{Expectations}\label{expectations}

\subsection{Required}\label{required}

\begin{itemize}
\tightlist
\item
  Do the readings and take part in discussion (30\%)
\item
  Take part in field experiment I (30\%)
\item
  Design an experiment as a candidate for survey experiment II (40\%)
\end{itemize}

\subsection{Optional}\label{optional}

\begin{itemize}
\tightlist
\item
  Prepare a research design or short paper, perhaps building on existing
  work. Typically this contains:

  \begin{itemize}
  \tightlist
  \item
    a problem statement
  \item
    a description of a method to address the problem
  \item
    analytic or simulation based results describing properties of the
    solution
  \item
    a discussion of implications for practice.
  \item
    A passing paper will illustrate subtle features of a method; a good
    paper will identify unknown properties of a method; an excellent
    paper will develop a new method.
  \end{itemize}
\end{itemize}

\section{Sessions and readings}\label{sessions-and-readings}

The course combines lectures, discussion, and activities.

All readings are linked from this document.

I will draw material especially from Blair, Coppock, and Humphreys
(2023): \href{https://book.declaredesign.org/}{Research Design in the
Social Sciences}

Recommended but not open access:
\href{https://wwnorton.com/books/9780393979954}{Gerber and Green (2012)}
with supplementary material available here:
\url{https://isps.yale.edu/FEDAI}

Readings for each session are given below.

\subsection{Intro}\label{intro}

Aims:

\begin{itemize}
\tightlist
\item
  Introduce the class
\item
  Practice running experiments by hand
\item
  Introduce and discuss field \href{in_class_experiment.html}{experiment
  I}
\item
  To prepare for next session: proposed adaptation of field experiment 1
\end{itemize}

\textbf{Course outline, tools}

\subsection{Experimentation}\label{experimentation}

Aims:

\begin{itemize}
\tightlist
\item
  Review lifecycle of an experiment (discussion)
\item
  Discuss ethical challenges to experimentation
\item
  Discuss a set of strong experiments (all graduate student projects)
\item
  Develop group field experiment 1
\end{itemize}

\textbf{To do before class}: 1. the readings; 2. develop specific
proposals in small groups for field experiment 1, completing
individually or collectively a copy of the
\href{https://egap.github.io/learningdays-resources/Exercises/design-form.html}{EGAP
research design form}

Required reading:

\begin{itemize}
\tightlist
\item
  \href{https://onlinelibrary.wiley.com/doi/full/10.1111/ajps.12627}{The
  Hijab penalty} and the
  \href{https://onlinelibrary.wiley.com/action/downloadSupplement?doi=10.1111\%2Fajps.12627&file=ajps12627-sup-0001-Appendix.pdf}{Appendix}
\item
  \href{https://www.cambridge.org/core/journals/ps-political-science-and-politics/article/protecting-the-community-lessons-from-the-montana-flyer-project/2595ED0DB8189897DF6BDB121BBB83CA}{Johnson:
  Protecting the Community: Lessons from the Montana Flyer Project}
\item
  \href{https://book.declaredesign.org/lifecycle}{Design lifecycle}
  Blair, Coppock, and Humphreys (2023) p 319 - 358
\item
  Example 1: Gerber, Green, and Larimer (2008)
\item
  Example 2: Lowe (2021)
\end{itemize}

Optional reading:

\begin{itemize}
\tightlist
\item
  \href{https://macartan.github.io/assets/pdf/papers/2015_reflections_ethics.pdf}{Humphreys
  (2015)} on ethics in experiments
\item
  \href{https://book.declaredesign.org/introduction/what-is-a-research-design.html}{DD,
  Ch 2} introduces ideas at a high level
\item
  \href{https://www.researchgate.net/publication/324240650_Research_Replication_Practical_Considerations/link/5ca3688b458515f7851d7b2e/download?}{Alvarez,
  Key, and Núñez (2018)}
\item
  \href{https://macartan.github.io/assets/pdf/papers/2013_fishing.pdf}{Humphreys,
  De la Sierra, and Van der Windt (2013)} on fishing and registration
\item
  Kurzban, Tooby, and Cosmides (2001)
\item
  Mousa (2020)
\end{itemize}

\subsection{Causality}\label{causality}

Aims:

\begin{itemize}
\tightlist
\item
  Understand the potential outcomes framework
\item
  Understand identification challenges and the role of randomization
\item
  Understand varieties of estimands
\end{itemize}

\textbf{To do before class}: 1. the readings; 2. collectively complete a
pre-analysis plan and ethics application for Experiment 1

Required:

\begin{itemize}
\tightlist
\item
  \href{https://www.jstor.org/stable/2289064}{Holland} is a beautiful
  classic reading on the fundamental problem of causal inference
\item
  \href{https://book.declaredesign.org/declaration-diagnosis-redesign/defining-inquiry.html}{DD,
  Ch 7}
\end{itemize}

Optional:

\begin{itemize}
\tightlist
\item
  \href{https://integrated-inferences.github.io/book/02-causal-models.html}{II
  Ch 2}, goes over both potential outcomes and DAGs
\item
  \href{https://hrr.w.uib.no/files/2019/01/Keele_2015_causal.pdf}{Keele
  (2015)} gives a good high level overview of many key ideas in this
  course
\item
  \href{https://gking.harvard.edu/files/matchse.pdf}{Imai, King, and
  Stuart (2008)}
\item
  \href{https://miguelhernan.org/whatifbook}{Hernán and Robins (2024)}
  (Section 3.1, 7.2, 8.4, 10.1)
\item
  \href{https://miguelhernan.org/whatifbook}{Hernán and Robins (2024)}
  Ch 6
\item
  \href{https://macartan.github.io/integrated_inferences/HJC4.html}{II
  Ch 4}
\item
  Lundberg, Johnson, and Stewart (2021)
\end{itemize}

\subsection{Answer strategies}\label{answer-strategies}

Aims:

\begin{itemize}
\tightlist
\item
  Understand alternative strategies to estimating treatment effects
\item
  Understand design based strategies for estimating standard errors and
  \(p\)-values
\item
  Understand confidence intervals
\item
  Understand the role of covariates in experimental design
\end{itemize}

\textbf{To do before class}: 1. Collective presentation on the
implementation and findings from Experiment 1

Optional:

\begin{itemize}
\tightlist
\item
  \href{https://www.stat.berkeley.edu/~census/neyregr.pdf}{Freedman
  (2008)} helps make connections between design based inference and
  regression
\item
  \href{https://scholar.harvard.edu/files/marie-abele/files/pnas.pdf}{Bind
  and Rubin (2020)} on RI
\item
  \href{https://arxiv.org/pdf/1208.2301.pdf}{Lin (2012)} relieves some
  worries you might have after reading Freedman (2008)
\item
  \href{https://onlinelibrary.wiley.com/doi/full/10.3982/ECTA20289}{Chang,
  Middleton, and Aronow (2024)} follows up on Lin (2012).
\end{itemize}

\subsection{Experimental Design and
Evaluation}\label{experimental-design-and-evaluation}

Aims:

\begin{itemize}
\tightlist
\item
  Understand alternative assignment strategies
\item
  Understand statistical power
\item
  Understand how to assess alternative diagnosands
\end{itemize}

\textbf{To do before class}: 1. The readings 2. Share initial design for
for survey based experiment.

Required:

\begin{itemize}
\tightlist
\item
  \href{https://book.declaredesign.org/declaration-diagnosis-redesign/diagnosing-designs.html}{DD,
  Ch 7}
\item
  \href{https://book.declaredesign.org/declaration-diagnosis-redesign/crafting-data-strategy.html}{DD,
  Ch 8} on data strategies
\item
  \href{https://book.declaredesign.org/declaration-diagnosis-redesign/choosing-answer-strategy.html}{DD,
  Ch 9} on answer strategies
\item
  \href{https://book.declaredesign.org/declaration-diagnosis-redesign/declaration-in-code.html}{DD,
  Ch 13} gives more practical guidance on getting started
\end{itemize}

\subsection{Topics}\label{topics}

Aims: * Topics TBD

\textbf{To do before class}: 1. The readings 2. Round-robin report on
design performance. Teams receive a class-mate's design, declare a model
that produces possible data from the design (without consulting the
authors design declaration, if any) and applies the original authors
answer srtategy, reporting performance.

Required reading (TBD):

\subsubsection{downstream
experimentation}\label{downstream-experimentation}

\begin{itemize}
\tightlist
\item
  Coppock and Green (2016)
\end{itemize}

\subsubsection{spillover experiments}\label{spillover-experiments}

\begin{itemize}
\tightlist
\item
  Sinclair, McConnell, and Green (2012)
\item
  \href{https://arxiv.org/abs/1305.6156}{Aronow and Samii (2017)} on
  spillovers
\end{itemize}

\subsubsection{patient preference
trials}\label{patient-preference-trials}

\begin{itemize}
\tightlist
\item
  Benedictis-Kessner et al. (2019)
\end{itemize}

\subsubsection{survey experiments}\label{survey-experiments}

\begin{itemize}
\tightlist
\item
  Hainmueller, Hopkins, and Yamamoto (2014)
\end{itemize}

\subsubsection{mediation experiments}\label{mediation-experiments}

\begin{itemize}
\tightlist
\item
  \href{https://imai.fas.harvard.edu/research/files/BaronKenny.pdf}{Imai,
  Keele, and Tingley (2010)} on mediation
\end{itemize}

\newpage

\section{Prerequisites and tools}\label{prerequisites-and-tools}

You should already have background in statistics up to the point of
feeling comfortable with regression.

In addition you should know some \texttt{R}. Really, the more you can
invest on getting on top of \texttt{R} before the class the better.

Before the first class please make sure your \texttt{R} is up-to-date
and that you are working in \texttt{R\ studio}. Then make sure you have
the following packages installed.

\begin{Shaded}
\begin{Highlighting}[]
\NormalTok{pacman}\SpecialCharTok{::}\FunctionTok{p\_load}\NormalTok{(}
\NormalTok{  rstan, }
\NormalTok{  dagitty, }
\NormalTok{  DeclareDesign, }
\NormalTok{  CausalQueries,}
\NormalTok{  tidyverse,}
\NormalTok{  knitr}
\NormalTok{  )}
\end{Highlighting}
\end{Shaded}

Resources:

\begin{itemize}
\tightlist
\item
  Resources for learning \texttt{R}:
  http://www.r-bloggers.com/how-to-learn-r-2/
\item
  Hadley Wickham's \href{https://r4ds.hadley.nz/}{R for Data Science}
\item
  \href{https://grantmcdermott.com/ds4e}{Grant McDermott on data
  science}
\end{itemize}

\subsection{\texorpdfstring{Writing with
\texttt{qmd}}{Writing with qmd}}\label{writing-with-qmd}

Please plan to do drafting in \texttt{Quarto}. This is a simple markup
language that lets you integrate writing and coding. This document is
written in \texttt{Quarto} and the slides will be also.

The key thing is that you can insert code chunks like this.

\begin{Shaded}
\begin{Highlighting}[]
\CommentTok{\# Define a random number}
\NormalTok{x }\OtherTok{\textless{}{-}} \FunctionTok{rnorm}\NormalTok{(}\DecValTok{1}\NormalTok{)}
\end{Highlighting}
\end{Shaded}

Code like this is run as the document compiles, and results can be
accessed as needed, like this: we just sampled the random number \(x =\)
2.1208917.

\begin{itemize}
\tightlist
\item
  See \href{https://r4ds.hadley.nz/}{Wickham, Çetinkaya-Rundel, and
  Grolemund (2023)} on
  \href{https://r4ds.hadley.nz/workflow-scripts}{Workflow} and
  \href{https://r4ds.hadley.nz/quarto}{Quarto}
\end{itemize}

I recommend using \texttt{Rstudio} as an editor. More information here:
https://quarto.org/docs/get-started/hello/rstudio.html

\subsection{AI}\label{ai}

AI can be a great aid to help understand this material and also to
develop code. However it should only be used as an aid: it makes many
many mistakes; you remain responsible for your code and more importantly
for your comprehension.

\section*{References}\label{references}
\addcontentsline{toc}{section}{References}

\phantomsection\label{refs}
\begin{CSLReferences}{1}{0}
\bibitem[\citeproctext]{ref-alvarez2018research}
Alvarez, R Michael, Ellen M Key, and Lucas Núñez. 2018. {``Research
Replication: Practical Considerations.''} \emph{PS: Political Science \&
Politics} 51 (2): 422--26.

\bibitem[\citeproctext]{ref-aronow2017estimating}
Aronow, P M, and Cyrus Samii. 2017. {``Estimating Average Causal Effects
Under General Interference, with Application to a Social Network
Experiment.''} \emph{Annals of Applied Statistics} 11 (4): 1912--47.

\bibitem[\citeproctext]{ref-de2019persuading}
Benedictis-Kessner, Justin de, Matthew A Baum, Adam J Berinsky, and
Teppei Yamamoto. 2019. {``Persuading the Enemy: Estimating the
Persuasive Effects of Partisan Media with the Preference-Incorporating
Choice and Assignment Design.''} \emph{American Political Science
Review} 113 (4): 902--16.

\bibitem[\citeproctext]{ref-bind2020possible}
Bind, M-AC, and DB Rubin. 2020. {``When Possible, Report a Fisher-Exact
p Value and Display Its Underlying Null Randomization Distribution.''}
\emph{Proceedings of the National Academy of Sciences} 117 (32):
19151--58.

\bibitem[\citeproctext]{ref-blair2023research}
Blair, Graeme, Alexander Coppock, and Macartan Humphreys. 2023.
\emph{Research Design in the Social Sciences: Declaration, Diagnosis,
and Redesign}. Princeton University Press.

\bibitem[\citeproctext]{ref-chang2024exact}
Chang, Haoge, Joel A Middleton, and PM Aronow. 2024. {``Exact Bias
Correction for Linear Adjustment of Randomized Controlled Trials.''}
\emph{Econometrica} 92 (5): 1503--19.

\bibitem[\citeproctext]{ref-coppock2016voting}
Coppock, Alexander, and Donald P Green. 2016. {``Is Voting Habit
Forming? New Evidence from Experiments and Regression
Discontinuities.''} \emph{American Journal of Political Science} 60 (4):
1044--62.

\bibitem[\citeproctext]{ref-freedman2008regression}
Freedman, David A. 2008. {``On Regression Adjustments to Experimental
Data.''} \emph{Advances in Applied Mathematics} 40 (2): 180--93.

\bibitem[\citeproctext]{ref-gerber2012field}
Gerber, Alan S, and Donald P Green. 2012. \emph{Field Experiments:
Design, Analysis, and Interpretation}. Norton.

\bibitem[\citeproctext]{ref-gerber2008social}
Gerber, Alan S, Donald P Green, and Christopher W Larimer. 2008.
{``Social Pressure and Voter Turnout: Evidence from a Large-Scale Field
Experiment.''} \emph{American Political Science Review} 102 (1): 33--48.

\bibitem[\citeproctext]{ref-hainmueller2014causal}
Hainmueller, Jens, Daniel J Hopkins, and Teppei Yamamoto. 2014.
{``Causal Inference in Conjoint Analysis: Understanding Multidimensional
Choices via Stated Preference Experiments.''} \emph{Political Analysis}
22 (1): 1--30.

\bibitem[\citeproctext]{ref-hernanwhatif}
Hernán, Miguel A, and James M Robins. 2024. \emph{What If?} Boca Raton:
Chapman \& Hall/CRC.

\bibitem[\citeproctext]{ref-humphreys2015reflections}
Humphreys, Macartan. 2015. {``Reflections on the Ethics of Social
Experimentation.''} \emph{Journal of Globalization and Development} 6
(1): 87--112.

\bibitem[\citeproctext]{ref-humphreys2013fishing}
Humphreys, Macartan, Raul Sanchez De la Sierra, and Peter Van der Windt.
2013. {``Fishing, Commitment, and Communication: A Proposal for
Comprehensive Nonbinding Research Registration.''} \emph{Political
Analysis} 21 (1): 1--20.

\bibitem[\citeproctext]{ref-imai2010general}
Imai, Kosuke, Luke Keele, and Dustin Tingley. 2010. {``A General
Approach to Causal Mediation Analysis.''} \emph{Psychological Methods}
15 (4): 309--34.

\bibitem[\citeproctext]{ref-imai2008misunderstandings}
Imai, Kosuke, Gary King, and Elizabeth A Stuart. 2008.
{``Misunderstandings Between Experimentalists and Observationalists
about Causal Inference.''} \emph{Journal of the Royal Statistical
Society Series A: Statistics in Society} 171 (2): 481--502.

\bibitem[\citeproctext]{ref-keele2015statistics}
Keele, Luke. 2015. {``The Statistics of Causal Inference: A View from
Political Methodology.''} \emph{Political Analysis} 23 (3): 313--35.

\bibitem[\citeproctext]{ref-kurzban2001can}
Kurzban, Robert, John Tooby, and Leda Cosmides. 2001. {``Can Race Be
Erased? Coalitional Computation and Social Categorization.''}
\emph{Proceedings of the National Academy of Sciences} 98 (26):
15387--92.

\bibitem[\citeproctext]{ref-lin2012agnostic}
Lin, Winston. 2012. {``Agnostic Notes on Regression Adjustments to
Experimental Data: Reexamining Freedman's Critique.''} \emph{arXiv
Preprint arXiv:1208.2301}.

\bibitem[\citeproctext]{ref-lowe2021types}
Lowe, Matt. 2021. {``Types of Contact: A Field Experiment on
Collaborative and Adversarial Caste Integration.''} \emph{American
Economic Review} 111 (6): 1807--44.

\bibitem[\citeproctext]{ref-lundberg2021your}
Lundberg, Ian, Rebecca Johnson, and Brandon M Stewart. 2021. {``What Is
Your Estimand? Defining the Target Quantity Connects Statistical
Evidence to Theory.''} \emph{American Sociological Review} 86 (3):
532--65.

\bibitem[\citeproctext]{ref-mousa2020building}
Mousa, Salma. 2020. {``Building Social Cohesion Between Christians and
Muslims Through Soccer in Post-ISIS Iraq.''} \emph{Science} 369 (6505):
866--70.

\bibitem[\citeproctext]{ref-sinclair2012detecting}
Sinclair, Betsy, Margaret McConnell, and Donald P Green. 2012.
{``Detecting Spillover Effects: Design and Analysis of Multilevel
Experiments.''} \emph{American Journal of Political Science} 56 (4):
1055--69.

\bibitem[\citeproctext]{ref-wickham2023r}
Wickham, Hadley, Mine Çetinkaya-Rundel, and Garrett Grolemund. 2023.
\emph{R for Data Science}. {O'Reilly Media, Inc.}

\end{CSLReferences}




\end{document}
